\documentclass[times, review, 10pt]{elsarticle}

\usepackage[top=4.3cm, right=4.8cm, bottom=4.3cm, left=4.8cm]{geometry}
\usepackage{hyperref}
\usepackage{booktabs}
\usepackage{tabularx}
\usepackage{multirow}
\usepackage{subcaption}
\usepackage{xcolor}
\usepackage{colortbl}
\usepackage{float}
\usepackage{graphicx}
\usepackage{makecell}
\usepackage{epstopdf}
\journal{Pattern Recognition}

% Override review mode line spacing to double spacing (2.0)
\renewcommand{\baselinestretch}{2.0}

\begin{document}

\begin{frontmatter}

\title{Interpretable supervised feature extraction for time series and spatial data}

\affiliation[1]{organization={Institute of Digitalization and General Medicine,
                Medical Faculty, RWTH Aachen University},
                addressline={Pauwelsstraße 30}, 
                city={Aachen},
                postcode={52074}, 
                country={Germany}}

\affiliation[2]{organization={Institute of Computational Biomedicine,
                Medical Faculty, RWTH Aachen University},
                addressline={Pauwelsstraße 30}, 
                postcode={52074}, 
                city={Aachen},
                country={Germany}}
                
\affiliation[3]{organization={Center for Rare Diseases Aachen,
                Medical Faculty, RWTH Aachen University},
                addressline={Pauwelsstraße 30}, 
                city={Aachen},
                postcode={52074}, 
                country={Germany}}

\affiliation[4]{organization={Department of Medical Statistics,
                Georg-August-University Göttingen},
                addressline={Humboldtallee 32}, 
                city={Göttingen},
                postcode={37073}, 
                country={Germany}}


\author[1,2]{J.C. Wolber}
\ead[cor1]{jwolber@ukaachen.de}
\cortext[cor1]{Corresponding author}

\author[4]{N.B. Paul}

\author[1,3]{J. Sellin}

\author[2]{M. E. Samadi}

\author[1,3]{M. Muecke}

\author[2]{A. Schuppert}

%% Abstract
\begin{abstract}
\textbf{Introduction}
Biomedical time series analysis relies on complex models that sacrifice interpretability for performance, limiting clinical adoption. 
Current automated methods like MiniRocket produce thousands of opaque features, while interpretable
approaches use fixed feature sets that miss task-specific patterns. We introduce SublimeX, an open-source framework for
interpretable feature extraction from time series and spatial data that bridges this gap.
\textbf{Methods}
SublimeX performs supervised feature extraction: at each iteration, Bayesian optimization proposes a candidate segment-based feature, a model is trained on
the current feature set augmented with the candidate, and cross-validation performance guides feature selection. Each feature computes the mean over an
optimized segment of a signal transformation (raw, z-score, derivative, or FFT power). The search stops when additional features no longer improve model
performance, yielding compact, task-specific feature sets.
\textbf{Results}
Across seven biomedical benchmarks, SublimeX achieves performance comparable to state-of-the-art methods while using only 5--20
extracted features versus thousands for competing approaches: best RMSE on short-term glucose trend forecasting (AZT1D, 31.11 vs.\ 31.18
MiniRocket), best AUC on gene expression prediction (REMC, $0.931 \pm 0.023$, tied with MiniRocket), 98.9\% accuracy on arrhythmia
classification (MITBIH), $0.831 \pm 0.020$ AUC on ICU mortality (MIMIC-IV; tsfresh leads at
$0.839 \pm 0.018$), and 86.0\% on activity recognition (PAMAP2).
\textbf{Discussion}
Post-hoc SHAP analysis on two contrasting domains shows that
discovered features align with established domain knowledge:
TSS-proximal histone marks (H3K4me3, H3K36me3, H3K27me3) for
gene expression, and recent glucose trajectory plus insulin and
carbohydrate effects for glucose forecasting---enabling
clinicians to inspect and validate model reasoning. An ablation
study confirms the importance of multi-transform search and
iterative optimization. SublimeX is available as an open-source
Python package via PyPI and GitHub.
\end{abstract}

%%Graphical abstract
\begin{graphicalabstract}
\includegraphics[width=\textwidth,keepaspectratio]{images/graphical_abtract/graphical_abstract.eps}
\end{graphicalabstract} 

%%Research highlights
\begin{highlights}
\item SublimeX is a novel feature extraction framework for discovering minimal, traceable features for Machine Learning.
\item Features are searched within segments of different signal transformations and channels via supervised optimization.
\item SublimeX achieves competitive performance on a wide range of biomedical tasks.
\item The extracted features are highly traceable and ideal for applications where transparency is essential.
\item SublimeX is available as an open-source software and can be customized to specific domains.
\end{highlights}

%% Keywords
\begin{keyword}
Feature extraction \sep Time series classification \sep Interpretability
\sep Feature-based methods \sep Bayesian optimization

\end{keyword}

\end{frontmatter}


\section{Introduction}
Biomedical machine learning increasingly operates on time series and structured signals: continuous bedside monitoring (ECG, EEG, oxygenation, glucose),
longitudinal electronic health records, wearable inertial sensors, and spatially indexed molecular measurements such as epigenomic tracks.
These modalities are central to both clinical decision support and biomedical discovery \cite{morid2023time, patharkar2024predictive}.
However, in many biomedical settings the bottleneck is not only predictive accuracy, but also whether a model produces evidence that is actionable and
auditable by domain experts \cite{tonekaboni2019clinicians, rudin2019stop}. For clinical adoption, predictions should be interpretable: clinicians need to
understand which aspects of the signal drove the outcome and to verify that the model's reasoning aligns with domain knowledge. Feature traceability---the
ability to map each model input back to a well-defined part of the raw signal (e.g., a channel, time segment, and transformation)---supports error
analysis, regulatory scrutiny, and trust, and enables localization of causative relations to specific signal regions (e.g., the P-wave in arrhythmia
detection) \cite{pantelidis2025ecgxplain}. Current automated and deep learning approaches often fall short by producing
high-dimensional or opaque representations that cannot be traced to specific signal regions or validated against physiological hypotheses. To address these shortcomings, we
present \textbf{SublimeX} (Supervised bottom-up Localized multi-representative Feature eXtraction), an interpretable interval-based feature extraction framework designed around three core principles that address the
limitations of existing approaches:

\begin{enumerate}
\item \textbf{Minimal feature set}: Rather than computing thousands of pre-defined features and
filtering post-hoc, SublimeX discovers task-specific features iteratively, adding one feature
at a time until no further improvement is observed. This yields compact representations tailored to each task.

\item \textbf{Traceable features}: Each SublimeX feature is \emph{traceable}: it
maps back to a concrete channel, transformation, and segment (the mean over an
optimized interval of raw, z-score, derivative, or FFT power on that channel).
Using localized segments rather than global (full-signal) statistics allows
discriminative information in specific regions to be captured. This traceability yields interpretable, human-readable
descriptions valuable for clinical applications where traceability and
localization of causal relationships is essential.

\item \textbf{Flexible and extensible framework}: SublimeX's modular architecture allows users
to customize signal transformations, aggregation functions, and machine learning models,
enabling domain-specific adaptations while retaining the core optimization framework.
\end{enumerate}

We validate SublimeX on seven diverse real-world biomedical
tasks: arrhythmia classification on MIT-BIH \cite{MoodyMark1992,Goldberger2000}, emotion
classification from EEG (Emotions) \cite{bird2018study}, gene expression prediction from
histone modification profiles (REMC) \cite{Roadmap2015}, mortality prediction from MIMIC-IV
clinical time series \cite{Johnson2023MIMICIV}, glucose forecasting (AZT1D)
\cite{khamesian2025azt1d}, human activity recognition from IMU sensors (PAMAP2)
\cite{Reiss2012PAMAP2}, and voice pathology detection (SVD) \cite{svd_database}. Crucially,
the inclusion of epigenomic data (REMC), treated as sequential data along genomic position
(not 2D/3D spatial), illustrates that SublimeX generalizes beyond temporal series to any
sequential data. 
We compare our approach against established automated baselines under matched splits and metrics.
We provide an ablation study to assess the
importance of SublimeX's key design choices, and provide
post-hoc SHAP analyses on two contrasting domains---gene
expression prediction and glucose forecasting---to
interpret the discovered features.

\section{Related Work}

We review three main paradigms for time series modeling---manual feature engineering,
automated feature extraction (including shapelet-based methods), and deep learning---before
situating SublimeX within this landscape.

\subsection{Manual feature engineering}
Traditional biomedical time series modeling relies on domain-driven descriptors (e.g., ECG
morphology, EEG band power) offering clear meaning and targeted error analysis
\cite{woldaregay2019data, fulcher2018feature}. It requires domain knowledge and experimentation with the training data to identify and tune effective descriptors. 
However, manual engineering is labor-intensive
and may miss task-specific characteristics across heterogeneous cohorts \cite{miotto2018deep}.


\subsection{Automated feature engineering}
Automated libraries compute large feature sets at scale. tsfresh
\cite{christ2018tsfresh} and catch22 \cite{lubba2019catch22} provide statistical coverage with
varying degrees of interpretability, while hctsa and TSFEL \cite{fulcher2017hctsa,
barandas2020tsfel} offer even larger descriptor sets. ROCKET and MiniROCKET
\cite{dempster2020rocket, dempster2021minirocket} achieve state-of-the-art accuracy through
random convolutional kernels with minimal computational overhead. Alongside such generic
methods, task-specific AI tools exist for particular modalities (e.g., deep learning for ECG arrhythmia detection \cite{yildirim2018arrhythmia}),
which we do not compare here. Shapelet-based methods \cite{hills2014shapelet, guillaume2022random}, another subset of automated feature engineering, use discriminative subsequences for distance-based features but focus on subsequence matching rather than regional statistics. These approaches compute fixed feature sets
regardless of task requirements: high-dimensional methods like MiniRocket (10,000 features)
sacrifice interpretability, while compact methods like catch22 (22 features) may miss
task-specific patterns not covered by their predefined descriptors.

\subsection{Deep learning approaches}
Deep learning methods learn representations end-to-end and often achieve benchmark-leading performance on time series tasks \cite{miotto2018deep, morid2023time}. Convolutional neural
networks (CNNs) have shown strong results in arrhythmia detection \cite{acharya2017ecg_deep, hannun2019cardiologist}, while recurrent architectures (LSTMs) excel at activity recognition
\cite{ordonez2016deep_activity, hammerla2016deep_activity}. 
These models discover complex temporal patterns without explicit feature engineering. However, their learned representations are difficult to interpret, 
validate against domain knowledge, or audit for failure modes---critical limitations in clinical settings where transparency is essential \cite{rudin2019stop, tonekaboni2019clinicians, wang2022systematic}.

\subsection{Comparison of time series classification approaches}
Table~\ref{tab:feature_toolkits} compares representative methods
across dimensions relevant to biomedical applications. We
select methods covering each major category in the taxonomy of
Middlehurst et al.\ \cite{middlehurst2024bakeoff} and the aeon toolkit \cite{middlehurst2024aeon}: interval-based,
feature-based, convolution-based, shapelet-based, dictionary-based,
deep learning, and hybrid approaches. SublimeX falls into the
interval-based category: like Time Series Forest (TSF)
\cite{deng2013tsf} and Canonical Interval Forest (CIF)
\cite{middlehurst2020cif}, it selects intervals of the signal and
computes summary statistics over them. However, unlike existing
interval-based methods that select intervals randomly, SublimeX
uses supervised Bayesian optimization to discover discriminative
intervals across multiple signal transformations.

From Table~\ref{tab:feature_toolkits}, we
select five baselines for our experimental evaluation
(Section~\ref{sec:baselines}). The selection balances
state-of-the-art performance, widespread adoption, and paradigm
diversity: tsfresh and catch22 represent widely used feature
engineering libraries, MiniROCKET is the leading convolution-based
method, RDST covers the shapelet paradigm, and a 1D CNN provides
an end-to-end deep learning baseline.

Our approach represents a unique combination of iterative feature discovery for time series and sequential data: we add one interpretable segment-based feature 
per round via supervised optimization, rather than constructing or selecting from a fixed feature set. Related iterative or incremental 
paradigms exist in other settings. Model-based experimental analysis (MEXA) \cite{marquardt2004model} applies iterative, 
incremental refinement of model structure for mechanistic modeling in process systems, but targets model identification and experiment design, not feature discovery. 
In explainable feature learning, Vlahek and Mongus \cite{vlahek2021efficient} propose an efficient iterative approach to learn interpretable features. 
Iterative feature construction for inductive learning \cite{piramuthu2009iterative} builds new features from existing ones using quality measures, 
while incremental forward feature selection \cite{lee2008incremental} adds features one at a time from a predefined pool (e.g., for microarray data). 
Unlike these, SublimeX discovers segment-based features from raw or transformed signals by optimizing segment boundaries and transform choice in a single search space, 
yielding traceable, task-specific features without a fixed feature library.

\input{tables/feature_toolkits.tex}

\section{Methods}

\subsection{Overview}
SublimeX discovers interpretable features from time series and sequential
data (including epigenomic tracks) through a supervised, iterative process
illustrated in Figure~\ref{fig:sublime_methodology}.  The procedure
has three stages.  First, a set of signal transformations is
precomputed once per input channel (\emph{Precompute
Transformations}).  Second, \emph{Iterative Feature Extraction} adds one
feature at a time.  In each round, a Bayesian optimizer runs 300~trials;
for every trial it calls an \emph{objective function} that (1)~receives
proposed feature parameters (channel, transform, segment boundaries),
(2)~computes the feature (extract segment, take mean → one scalar per
sample), and (3)~evaluates the model on all current features plus the
candidate and returns the validation score.  The optimizer uses these
scores to propose the next trial; after 300~trials the best candidate is
added to the feature set and a new round starts.  The loop stops when no
candidate improves the score (\emph{stagnation}).  Third, the resulting
compact feature matrix is returned as the \emph{Final Output}.

\begin{figure}[H]
  \centering
  \includegraphics[width=0.65\textwidth,keepaspectratio]{images/methods/flowchart.png}
  \caption{SublimeX methodology overview.
  \textbf{Input Data}: time series or sequential data
  (univariate or multivariate).
  \textbf{Precompute Transformations}: four signal
  transformations (raw, z-score, derivative, FFT power)
  are applied once per channel.
  \textbf{Iterative Feature Extraction}: one feature is
  added per round.  Each round starts with the current
  feature set; the Bayesian optimizer (TPE) calls the
  \emph{objective function} (dashed box) 300~times.  The
  objective function \textbf{(1)}~receives proposed
  feature parameters (channel, transform, segment
  boundaries), \textbf{(2)}~computes the feature
  (extract segment, mean → one scalar per sample), and
  \textbf{(3)}~evaluates the model on existing features
  plus the candidate and returns the validation score.
  After 300~trials the best candidate is added to the
  feature set (\emph{improvement}) and the next round
  starts; on \emph{stagnation} (no improvement) the loop
  stops and the \textbf{Final Output} is the compact
  feature matrix.}
  \label{fig:sublime_methodology}
\end{figure}

\subsection{Signal Transformation}
The first stage applies a set of signal transformations to all input channels, creating
multiple views of the data. The default transformation library includes four transforms:
\textbf{Raw} (identity transform preserving the original signal values), \textbf{Z-score}
(per-sample standardization $z_t = (x_t - \mu) / \sigma$, removing amplitude differences to
focus on shape), \textbf{Derivative} (first-order gradient $\nabla x_t$, capturing rate of
change and highlighting transitions), and \textbf{FFT Power} (power spectrum $|F(x)|^2$
interpolated back to the original time length, capturing frequency content). Users can extend
the library with domain-specific transformations, enabling expert knowledge in the feature discovery process.

All transformations are precomputed once and stored, creating parallel views of the input signals. During feature optimization, the optimizer selects
which transformation to use for each feature and, for multivariate data, which channel; discriminative information can thus emerge in whichever
signal representation is most predictive for the task.

\subsection{Feature Extraction}
The core of SublimeX is a supervised optimization loop that discovers discriminative
segment-based features (Figure~\ref{fig:sublime_methodology}).  Each feature is defined
by \textbf{transform type} (which signal view: raw, z-score, derivative, or FFT power),
\textbf{segment boundaries} (start and end indices of a fixed interval), and, for
multivariate data, \textbf{channel index}.  The same segment is applied to every
sample; each sample yields one scalar per feature (e.g., the mean over that segment).

The optimizer searches over these parameters by repeatedly calling an \textbf{objective
function}.  One call corresponds to one candidate feature: the objective (1)~receives
proposed parameters, (2)~computes the feature (extract the segment from the chosen
transformed channel, take the mean), and (3)~evaluates the downstream model on the
current feature set plus the candidate, returning the validation score.  The Bayesian
optimizer (TPE) uses these scores to propose the next candidate; after 300~trials per
round, the best candidate is added to the feature set and a new round begins.  The
loop stops when no candidate improves the score (stagnation), yielding a compact,
interpretable feature matrix.

For univariate series the channel is fixed; for multivariate series the optimizer
chooses the channel per feature.  The procedure is greedy (one feature per round) and
may converge to local optima; every retained feature is kept in the feature set for
subsequent rounds. 

\subsection{Implementation Details}
Our default configuration is as follows. We use Bayesian optimization with Optuna's Tree-structured Parzen Estimator (TPE)
sampler \cite{bergstra2011tpe, watanabe2023tpe} to search the space of feature
parameters (channel, transform, segment start and end), with 300 trials per feature.
TPE is well suited to this setting because the search space is complex and mixed: it combines categorical variables (channel index, transform type) with continuous variables (segment boundaries), which standard gradient-based or grid-search methods handle poorly; TPE models the conditional distribution of good versus poor trials and efficiently explores such mixed spaces \cite{bergstra2011tpe}.
The procedure is analogous to automated hyperparameter optimization, with feature parameters (channel, transform, segment boundaries) as the search variables and validation performance as the objective \cite{bergstra2011tpe}. Bayesian optimization with TPE is one of several potentially useful optimizers; the framework is agnostic to the choice (e.g., we also evaluate NSGA-II in the ablation study).

Internal validation uses a single deterministic 50/50 train/validation split (fixed
random seed) to evaluate candidate features during optimization \cite{bergstra2011tpe};
every trial is evaluated on the same held-out
samples, ensuring that performance comparisons are not confounded by split
variability. The default transformation library does not include wavelets (only raw,
z-score, derivative, and FFT power). The default machine learning model is LightGBM
with maximum tree depth of 5.

A key strength is SublimeX's modular design: users can customize pipeline components for domain-specific needs. 
Specifically, SublimeX accepts custom objective functions (e.g. one may not want to use the mean aggregation function but instead the median or the maximum), custom ML models (any
scikit-learn-compatible model with fit/predict interface; default: LightGBM; alternatives include linear models, SVMs, or neural networks), and custom transforms 
(domain-specific signal transformations beyond the default library). 
This lets practitioners incorporate domain knowledge while retaining the core optimization loop. 
SublimeX is available as an official Python package on PyPI (see \url{https://pypi.org/project/sublimex/}). We ran our experiments on a MacBook Pro M2 using 12 cores and 32GB of RAM with Python 3.13.5.

\subsection{Datasets}
We evaluate SublimeX on seven biomedical datasets: glucose forecasting, arrhythmia classification, emotion recognition, gene expression prediction, clinical outcome prediction, 
activity recognition, and voice pathology detection. Table~\ref{tab:datasets} summarizes their characteristics. Preprocessing and data sources are in the Supplementary Material.
We select them for their diversity in tasks, domains, and signal characteristics.

\input{tables/dataset_characteristics.tex}

\subsection{Baselines}
\label{sec:baselines}
We compare SublimeX against five time series methods: tsfresh \cite{christ2018tsfresh}, catch22 \cite{lubba2019catch22}, MiniRocket \cite{dempster2021minirocket}, 
RDST \cite{guillaume2022random}, and a 1D CNN. For all feature-based methods (tsfresh, catch22, MiniRocket, RDST, SublimeX), we train a LightGBM
classifier \cite{ke2017lightgbm} on the extracted features using default hyperparameters. This ensures that performance differences reflect the quality of the extracted features rather than
classifier tuning. Implementation details are in the Supplementary Material.

\subsection{Evaluation}
We evaluate all methods using stratified 5-fold cross-validation for most datasets. For PAMAP2, we follow Tello et al.\ \cite{tello2025contrasting} protocol with 512-sample windows
(5.12 seconds) and 100-sample shift (1 second), combined with leave-one-subject-out (LOSO) cross-validation to ensure subject-independent evaluation, as recommended by recent work on
HAR evaluation protocols \cite{tello2025contrasting, geissler2024power}. For the AZT1D glucose forecasting dataset, we use a one time temporal 80\%/20\% train/test split per patient to prevent data leakage.
For REMC, we run experiments on all 56 cell lines from the Roadmap Epigenomics Consortium, evaluating each cell line independently with 5-fold CV.
Results are aggregated as mean $\pm$ standard deviation across folds for standard datasets, across all cell lines and folds for REMC, across
all subjects for PAMAP2 (8 LOSO folds), and across all subjects for AZT1D. Standard datasets use the same 5-fold CV protocol as the main evaluation. Evaluation metrics are task-dependent:
AUC for REMC and MIMIC, binary accuracy for SVD, multi-class classification accuracy for Emotions, MITBIH, PAMAP2 and root mean squared error (RMSE) for AZT1D.
Evaluation metrics are in Table~\ref{tab:datasets}.

\section{Results}

\subsection{Performance Comparison}
\input{tables/results_table.tex}

Table~\ref{tab:results} compares SublimeX against five baselines across seven biomedical datasets.
SublimeX performs competitively throughout and leads on two datasets.
On REMC, it matches MiniRocket with the highest AUC ($0.931 \pm 0.023$), and on AZT1D it obtains the lowest RMSE ($31.11$), ahead of MiniRocket ($31.18$) and RDST ($31.91$).

Performance remains close to top-scoring methods on the remaining datasets. On MIMIC-IV mortality prediction, tsfresh achieves the best AUC ($0.839 \pm 0.018$), with SublimeX
close behind ($0.831 \pm 0.020$) and MiniRocket and CNN at $0.825$--$0.826$; the small spread
indicates that multiple feature representations capture prognostic information in ICU time series.
On MITBIH, accuracy is $0.989 \pm 0.001$, narrowly behind MiniRocket ($0.992$) and RDST ($0.991$).
The gap is somewhat larger on Emotions ($0.942 \pm 0.014$ vs.\ MiniRocket's $0.988 \pm 0.007$) and PAMAP2 ($0.860 \pm 0.081$ vs.\ tsfresh's $0.898$ and catch22's $0.893$).
For SVD voice pathology detection, RDST leads at $0.718 \pm 0.016$ while SublimeX scores $0.689 \pm 0.014$; all methods remain below 72\% on this challenging task.

\subsection{Transform Composition and Pattern Stability}

Beyond predictive accuracy, SublimeX's learned features
can be inspected: each has a direct mathematical definition (channel, transform,
segment) and is traceable to the raw signal. This traceability is by design---users can
map a feature to a specific channel, transform, and segment---supporting
interpretability before domain validation. Medical interpretability
requires domain validation (e.g., SHAP analysis and alignment with prior knowledge, as
in Section~\ref{sec:domain_interpretation}). We investigated which signal
transforms and segments SublimeX selected, and
how stable these are across five cross-validation folds.
Feature discovery is not deterministic because each fold uses different samples.
Figure~\ref{fig:performance_analysis} summarizes the selected transformations and pattern stability across folds.

\begin{figure}[H]
  \centering
  \includegraphics[width=\textwidth,keepaspectratio]
  {images/results/performance_stability_analysis.png}
  \caption{Transform composition and pattern stability.
  (A)~Proportion of extracted features per transformation per dataset.
  (B)~Segment position of the \emph{first} extracted feature across folds.
  Bars show mean segment position; error bars show variability.
  T = transform consistency, C = channel consistency (fraction of folds with same transform/channel).}
  \label{fig:performance_analysis}
\end{figure}

Panel~A aggregates the transforms chosen across all features per dataset, revealing which representations SublimeX favors.
For MITBIH arrhythmia classification, derivatives (49\%) and z-score normalization (32\%) capture QRS morphology independent of amplitude. For AZT1D glucose forecasting, 
raw values (44\%) and derivatives (31\%) dominate, reflecting the importance of absolute glucose levels and rates of change for predicting future dynamics. 
For amplitude-driven tasks like PAMAP2 (IMU sensors) and REMC (ChIP-seq signal), raw values dominate (58--64\%). MIMIC shows balanced usage across all four transforms, 
while Emotions and SVD combine raw values with FFT power.

Panel~B focuses on the \emph{first} feature only (the one selected first in the
iterative search; selection order can vary across folds due to local optima), showing
its mean segment position across cross-validation folds with error bars indicating
positional spread. The first selected feature often captures global signal behavior;
later features tend to add more detail, though counterexamples exist.
MITBIH achieves strong stability (position std~=~0.027, T:0.8), with most folds identifying nearly identical cardiac-cycle regions. PAMAP2 shows low
positional variability (std~=~0.039) with high channel consistency (C:0.8). In contrast, REMC shows higher positional variability (std~=~0.31) with T:0.8,
reflecting multiple valid representations of histone modification patterns near the TSS. Emotions and SVD show moderate positional variability (std~=~0.23 and~0.21) despite perfect transform consistency, 
suggesting multiple discriminative regions and different local optima.
MIMIC has the lowest transform consistency (T:0.4), with multiple representations carrying complementary prognostic information.

\subsection{Incremental Feature Analysis}
\label{sec:feature_analysis}

To characterize why SublimeX succeeds on some datasets but not others, we conducted an incremental feature analysis of
train, validation, and test performance as features are added (Figure~\ref{fig:incremental_features}). We use the same
cross-validation splits as the main evaluation and reconstruct the
50/50 train/validation split used during feature optimization. A LightGBM model is trained on the optimization-train
subset and evaluated on all three splits.

\begin{figure}[H]
  \centering
  \includegraphics[width=\textwidth,keepaspectratio]
  {images/results/incremental_features.png}
  \caption{Incremental feature analysis showing train, validation,
  and test performance as features are added. The model is trained on
  the optimization-train subset (the same data used during feature
  search) and evaluated on all three splits. Metrics vary by dataset:
  accuracy (Emotions, MITBIH, SVD, PAMAP2), AUC (MIMIC, REMC), and
  RMSE (AZT1D; lower is better). Solid lines show mean performance
  across folds; shaded areas indicate $\pm$1 standard deviation.
  Note that across folds different amounts of features are extracted
  and thus standard deviation may not be available for all incremental
  evaluations.}
  \label{fig:incremental_features}
\end{figure}

MITBIH and REMC show tight train--validation--test convergence
(validation--test gaps ${\sim}0.1\%$ and ${\sim}0.5\%$),
indicating that discovered features capture
discriminative patterns in datasets with localized,
consistent signal structure (QRS morphology, TSS-proximal
histone enrichment).
Emotions and MIMIC show large train--test gaps
(${\sim}6\%$ and ${\sim}17\%$) driven by LightGBM overfitting
to the small training set, but their validation--test gaps
remain small (${\sim}2\%$ and ${\sim}1.6\%$), indicating
that the feature search itself generalizes well.
SVD shows moderate validation--test separation
(${\sim}6\%$), suggesting that the feature search partially
overfits to the optimization split on this challenging task.
PAMAP2 exhibits the most extreme pattern: training accuracy
reaches 100\% and validation 98.7\%, yet test accuracy
plateaus at ${\sim}73\%$ (validation--test gap
${\sim}26\%$). This large gap arises because the internal
50/50 train/validation split shares the same subjects,
whereas LOSO evaluation tests on an entirely unseen subject,
confirming that features capture
subject-specific movement patterns rather than
generalizable activity signatures.
For AZT1D, validation and test RMSE converge
to similar values (31.9 and 30.9), indicating robust feature
discovery for glucose prediction. Test RMSE is slightly lower than
validation RMSE, which is plausible under the temporal split: test
uses the last 20\% of each patient's timeline while validation is
drawn from the first 80\%, so the test segment can exhibit different
predictability (e.g., lower target variance or more stable regimes)
without implying leakage.

\subsection{Ablation Study}

We ablate key design choices across seven datasets.
All ablation experiments are evaluated on a single fold
only (fold~1), so no standard deviations are reported and
Baseline scores may differ slightly from the multi-fold
averages in Table~\ref{tab:results}.
To further reduce computational cost, REMC experiments
use only the first cell line (E003) and PAMAP2 experiments
use only the first subject (single LOSO fold).
Our default configuration uses mean aggregation, 300 trials
per feature, a single 50/50 train/validation split for
internal evaluation, the full transform library (raw,
z-score, derivative, FFT power), and the TPE sampler.
Table~\ref{tab:ablation} reports score, feature count,
and runtime (in seconds) for each variant.

We evaluate seven variants against the Baseline. Baseline serves as our default configuration using mean aggregation with LightGBM and the TPE sampler.
\emph{Optimized Aggregate} expands the search space to include eight aggregation functions (mean, min, max, range, std, median, argmin, argmax). 
\emph{Pattern Search} searches for discriminative quadratic B\'ezier patterns using Euclidean distance similarity, similar to shapelet-based methods such as RDST \cite{guillaume2022random}.
\emph{Decision Tree} replaces LightGBM with a Decision Tree (max\_depth=5) during optimization and test set prediction, testing whether a simpler internal model improves feature quality. 
\emph{1000 Trials} increases the budget from 300 to 1000 trials per feature.
\emph{Raw Only} disables signal transforms, using only raw signals. \emph{NSGA-II} replaces TPE with a multi-objective evolutionary algorithm. 
\emph{Parallel} extracts all features simultaneously in a single optimization run (using the same number of features as Baseline) rather than iteratively, 
testing whether joint optimization outperforms greedy sequential search.

\begin{table}[H]
\centering
\caption{Ablation study across seven datasets (single fold;
Baseline scores may differ slightly from the multi-fold
averages in Table~\ref{tab:results}). Bold indicates best
score per dataset. Arrows indicate optimization direction
($\uparrow$ higher better, $\downarrow$ lower better).
Rows show metric, number of extracted features, and
runtime (s).}
\input{tables/ablation_results.tex}
\end{table}

Signal transforms are
critical: the Raw Only variant shows degradation on Emotions
(93.2\% vs.\ 94.4\% Parallel) and MITBIH (98.3\% vs.\
98.9\% NSGA-II), confirming that multi-representation
search across raw, z-score, derivative, and FFT power
transforms is essential for capturing discriminative
information. Despite its simplicity, mean aggregation is
effective: the Baseline achieves best or near-best
performance on Emotions, REMC, MITBIH, and SVD,
demonstrating that regional mean statistics capture
discriminative information effectively.
On AZT1D, NSGA-II achieves the best RMSE (30.81),
closely followed by the Aggregate variant (30.88),
while Raw Only scores worst (31.25), underscoring
the value of multi-transform search for glucose
forecasting.
The Decision Tree variant, which uses a depth-5 decision
tree for both optimization and final prediction, shows
substantial performance degradation compared to
LightGBM-based variants (e.g., 96.2\% vs.\ 98.9\% on
MITBIH, 0.785 vs.\ 0.855 Baseline on MIMIC-IV),
confirming that model capacity matters for test-time
evaluation even when feature discovery is identical.

On MIMIC-IV, the Parallel variant achieves the highest single-fold AUC (0.857), slightly ahead of the
Baseline (0.855) and Raw Only (0.850), indicating that for ICU mortality prediction joint optimization
of multiple features can marginally outperform the default iterative greedy search, while the balanced
use of transforms (Baseline) remains highly competitive. 
Additional trials show diminishing returns: the 1000 Trials variant does not consistently outperform the Baseline (300 trials), indicating
that the default budget is sufficient for most tasks. NSGA-II performs
comparably to TPE across datasets, with tied best performance on MITBIH
(0.989). The Parallel variant achieves strong performance on Emotions (94.4\%
vs.\ 94.8\% Baseline on fold 1), demonstrating that joint optimization of all
features can approach greedy sequential search performance. However, Parallel
shows mixed results across datasets: it matches Baseline feature counts exactly
(by design) but achieves lower performance on MITBIH (98.6\% vs.\ 98.9\% on
fold 1), REMC (0.926 vs.\ 0.928 on fold 1), and SVD (66.8\% vs.\ 69.3\% on
fold 1), suggesting that the iterative greedy strategy captures complementary
information that joint optimization misses. Notably, Parallel is substantially
faster (5--77 seconds vs.\ hundreds or thousands for iterative variants),
making it attractive for rapid prototyping despite the performance trade-off.

\subsection{Domain-Specific Feature Interpretation}
\label{sec:domain_interpretation}

To demonstrate that SublimeX's features admit meaningful
domain interpretation, we apply SHAP \cite{lundberg2017unified} to two tasks
with distinct clinical contexts: REMC gene expression and AZT1D glucose
trend forecasting. Figure~\ref{fig:domain_interpretation} shows the
visualizations and verifiable results (panels A--F in the caption).

\subsubsection{REMC: Gene Expression Prediction from
Histone Modifications}

In molecular biology, histone modifications are chemical
tags attached to histone proteins---the ``spools'' around
which DNA is wound. These tags act as molecular switches
that control whether nearby genes are turned on (expressed)
or off (silenced), without altering the DNA sequence itself.
Five common modifications---H3K4me3, H3K4me1, H3K36me3,
H3K9me3, and H3K27me3---form a well-characterized
regulatory code: certain marks (e.g.\ H3K4me3) are
consistently enriched near the start of active genes,
while others (e.g.\ H3K27me3) accumulate at silenced loci
\cite{karlic2010histone, howe2017h3k4me3}.

Using the 12~SublimeX features extracted for cell line E003
(fold~1), we train LightGBM on the extracted features alone
(AUC~=~0.928) and compute SHAP values.
Panels~A and~B focus on the most important feature,
H3K4me3 raw 65--96---the mean H3K4me3 enrichment signal
in the region surrounding the transcription
start site (TSS), where genes are read. Panel~A
shows that actively expressed genes exhibit a strong
enrichment peak in this region, while silenced genes remain
flat. Panel~B confirms clear class separation (mean 16.9
for high expression vs.\ 4.0 for low expression).

Panel~C shows the SHAP beeswarm plot for all 12~features.
The most important feature is H3K4me3 raw 65--96
(mean absolute SHAP~=~1.47), followed by H3K36me3 raw
0--61 (mean absolute SHAP~=~0.86), which captures
H3K36me3 enrichment along the gene body---a hallmark of
actively transcribed genes. The repressive mark H3K27me3
ranks third (mean absolute SHAP~=~0.51), consistent
with its known role in gene silencing. This ranking
precisely mirrors the established epigenomic hierarchy
\cite{karlic2010histone, howe2017h3k4me3}: the activating
mark near the TSS is the strongest predictor, followed by
the elongation-associated mark along the gene body, and
then the silencing mark.

\subsubsection{AZT1D: Glucose Forecasting from
Continuous Monitoring Data}

Continuous glucose monitoring (CGM) devices measure blood
sugar levels every five minutes, producing a dense time
series that captures the dynamic interplay between meals,
insulin doses, physical activity, and metabolic state.
Short-term glucose trend prediction (60~minutes
ahead) is clinically important for people with diabetes,
enabling proactive adjustments to insulin and
carbohydrate intake to avoid dangerous hypo- or
hyperglycemic episodes.

Using the 16~SublimeX features extracted for AZT1D, we train LightGBM on the training set with the extracted features and compute SHAP values. Panels~D and~E
focus on the most important feature, CGM raw 16--23---the mean glucose level over the most recent 35 minutes of CGM history.
We group samples into ``rising'' vs.\ ``falling'' glucose by splitting at the
median of the 60-minute-ahead glucose change (the regression target), so that
we can compare class-conditional distributions despite the task being regression.
Panel~D shows that samples where glucose is about to rise (red) already have elevated recent glucose levels
(mean ${\sim}$164~mg/dL), while samples where glucose
is about to fall (blue) show higher recent readings
(mean ${\sim}$130~mg/dL) that are declining---a pattern
consistent with the post-meal glucose trajectory, where an
initial rise is followed by insulin-mediated decline.

Panel~F shows the SHAP beeswarm plot for all 16~features.
The dominant feature is CGM raw 16--23 (mean
absolute SHAP~=~11.77), confirming that the current
glucose trajectory is the primary predictor. The second
and third features---CGM deriv.\ 21--23 (mean
absolute SHAP~=~3.57) and CGM raw 22--23 (mean
absolute SHAP~=~2.94)---capture the instantaneous rate of
glucose change and the most recent glucose reading,
respectively. Notably, the fourth- and fifth-ranked features
are Insulin FFT 0--10 (mean absolute SHAP~=~1.60) and Carbs
deriv.\ 10--23 (mean absolute SHAP~=~1.41), reflecting the
frequency pattern of insulin dosing and the rate of
carbohydrate absorption effect, both influencing future glucose dynamics. This hierarchy
aligns with clinical intuition: current glucose level and
trend are the strongest short-term predictors; insulin
and carbohydrate effects provide secondary but clinically
essential context for anticipating glucose excursions
\cite{khamesian2025azt1d}.

\begin{figure}[H]
  \centering
  \includegraphics[width=\textwidth,keepaspectratio]
  {images/results/domain_interpretation.png}
  \caption{Domain interpretation via SHAP analysis for two biomedical tasks (verifiable
  results). \textbf{Top row---REMC gene expression} (cell line E003, fold~1):
  \textbf{(A)}~Mean H3K4me3 signal within the SublimeX-extracted segment (bins~65--96,
  near the TSS) for highly vs.\ lowly expressed genes;
  \textbf{(B)}~class-conditional density of that H3K4me3 feature;
  \textbf{(C)}~SHAP beeswarm for all 12~features (feature importance ranking).
  \textbf{Bottom row---AZT1D glucose trend prediction} (fold~1):
  \textbf{(D)}~Mean CGM signal in the top feature segment ($-$40 to 0~min) for
  samples with rising vs.\ falling glucose 60~min ahead;
  \textbf{(E)}~distribution of the top CGM feature for those two groups;
  \textbf{(F)}~SHAP beeswarm for all 16~features. In beeswarm plots, each dot is one
  test sample; color = feature value (red high, blue low), horizontal position = SHAP
  contribution.}
  \label{fig:domain_interpretation}
\end{figure}

The features themselves are mathematical (channel, transform, segment, mean);
medical interpretability arises from post-hoc validation: SHAP and alignment with
domain knowledge show that SublimeX autonomously discovers clinically and
biologically meaningful feature hierarchies---the epigenomic activation hierarchy for
gene expression and the glucose-trend-then-treatment hierarchy for diabetes
management---without any domain-specific guidance. Crucially, in both examples the
highest impact on model prediction can be traced to a very localized segment or region:
for REMC, the dominant feature is H3K4me3 in a narrow window near the TSS (bins 65--96),
allowing those genomic regions to be targeted for epigenetic discovery or
intervention; for AZT1D, the dominant feature is mean glucose over the most recent
35 minutes, so clinical interventions (e.g., insulin or carbohydrate adjustment)
can focus on that window. This localization opens possibilities for
scientific discovery and targeted clinical action.

\section{Discussion}

\subsection{Summary of Main Findings}

Across seven biomedical datasets, SublimeX achieves competitive
performance while using only 5--20 extracted features versus
thousands for competing approaches. It leads on AZT1D glucose
forecasting (31.11 RMSE) and ties with MiniRocket on REMC gene
expression ($0.931 \pm 0.023$ AUC across 56 cell lines), with
strong results elsewhere: $0.831 \pm 0.020$ AUC on MIMIC-IV (tsfresh leads at $0.839 \pm 0.018$),
98.9\% on MITBIH, 94.2\% on Emotions, and 86.0\% on PAMAP2.

Discovered features show task-dependent
transform selection---derivatives and z-score for MITBIH (QRS
morphology), raw values for amplitude-driven tasks (PAMAP2, REMC),
balanced usage for MIMIC---with high positional stability on
MITBIH and PAMAP2 and greater variability on REMC and MIMIC, where
multiple discriminative regions exist. Incremental feature analysis
(Section~\ref{sec:feature_analysis}) shows that SublimeX excels
when signal structure is localized and consistent: MITBIH and REMC
exhibit tight train--validation--test convergence (gaps
${\sim}0.1\%$ and ${\sim}0.5\%$), whereas PAMAP2 suffers a
validation--test gap of ${\sim}26\%$ because LOSO evaluation tests
on entirely unseen subjects whose movement patterns differ from
those seen during optimization.

Ablation experiments confirm that multi-transform search is
essential (Raw Only degrades Emotions and MITBIH), mean aggregation
is sufficient for most tasks, and the 300-trial default budget
shows diminishing returns when increased to 1000. The iterative
greedy search outperforms joint parallel optimization on several
datasets, indicating that sequential feature addition captures
complementary information. SHAP analysis on two contrasting
biomedical domains demonstrates that SublimeX autonomously
discovers meaningful feature hierarchies in both cases: the
established epigenomic hierarchy for REMC---H3K4me3 near the
TSS, H3K36me3 in the gene body, and the repressive mark H3K27me3
\cite{karlic2010histone, howe2017h3k4me3}---and the clinically
intuitive glucose-trend-then-treatment hierarchy for AZT1D, where
current glucose level and its rate of change dominate, followed by
insulin and carbohydrate effects
\cite{khamesian2025azt1d}. These results illustrate the
interpretability that compact, explicitly described features
afford across fundamentally different data modalities and clinical
questions.

\subsection{Comparison to Related Work}

SublimeX performance compares favorably with task-specific approaches reported in the literature. For MITBIH arrhythmia classification, deep learning
approaches typically achieve 94--99\% accuracy but require thousands of parameters and GPU inference \cite{acharya2017ecg_deep, hannun2019cardiologist}.
Our 98.9\% accuracy with interpretable features represents a practical configuration for clinical settings where model transparency and computational
simplicity are prioritized. For MIMIC mortality prediction, published AUCs range from 0.75 to 0.85 depending on the feature set and time horizon
\cite{harutyunyan2019multitask, purushotham2018benchmarking}. SublimeX's $0.831 \pm 0.020$ AUC (mean across 5 folds; fold 1: 0.855) falls within this
range while offering explicit, interpretable features that support clinical validation;
tsfresh reaches $0.839 \pm 0.018$ on the same task.

For AZT1D glucose forecasting, our global SublimeX model trained on all patients combined achieves
31.11~RMSE, which is higher than the ${\sim}25$~RMSE reported in our prior work using
personalized, patient-specific models that integrate mechanistic components with multimodal large
language models \cite{wolber5261056multimodal}. This difference is expected and can be partially
explained by the fact that, in this work, we train the global model only once on the initial
training set and evaluate on held-out samples, whereas prior work retrained the model anew for
every prediction interval (e.g., every hour), thereby leveraging the most up-to-date patient data.
Much of the glucose-forecasting literature similarly uses per-patient or dynamic retraining
approaches---or highly specialized pipelines---to minimize raw error
\cite{woldaregay2019data}. In contrast, our goal here is to keep the experimental setup simple,
enable a fair, like-for-like comparison across all approaches, and focus on the methodology and
interpretability of SublimeX rather than maximizing absolute glucose forecasting accuracy for
AZT1D.

On PAMAP2 activity recognition with LOSO evaluation, SublimeX achieves
$0.860 \pm 0.081$ accuracy, while tsfresh achieves $0.898 \pm 0.076$. This
result must be interpreted in light of recent methodological critiques: Tello
et al.\ \cite{tello2025contrasting} demonstrated that conventional HAR
evaluation practices---combining sliding windows with standard k-fold
cross-validation---systematically overestimate accuracy by allowing temporal
leakage between training and test sets. Many published PAMAP2 results,
including the original 77.8\% accuracy \cite{reiss2013confidence} and deep
learning approaches reporting 85--90\% \cite{ordonez2016deep_activity}, used
evaluation protocols that may inflate performance estimates. Our LOSO protocol
follows Tello et al.'s recommendations, ensuring subject-independent evaluation
without temporal leakage. Notably, all methods show substantial inter-subject
variability on this task: CNN achieves $0.790 \pm 0.180$, and RDST achieves
$0.819 \pm 0.187$, indicating that subject-dependent generalization remains challenging regardless of method.

Biomedical deployment places constraints beyond benchmark accuracy: clinicians and scientists must be able to \emph{inspect} what evidence
a model uses, relate it to domain knowledge, and understand when and why failures occur \cite{tonekaboni2019clinicians, caruana2015intelligible}.
SublimeX represents each sample through segment-based features with explicit, human-readable
descriptions, attractive for biomedical use because features are interpretable, localized to specific signal regions, and compact enough for
standard attribution methods (e.g., SHAP) to yield meaningful explanations \cite{fulcher2018feature}. Traceability to channel, transform, and segment pins the main drivers of a prediction to narrow regions---as in REMC (TSS-proximal histone marks) and AZT1D (recent glucose window)---enabling knowledge discovery and targeted intervention (e.g., epigenetic targeting of loci or timing of glucose-related actions).

\subsection{Limitations}

The iterative addition of features and the predefined transform library can also be limitations. Runtime can be
substantial when increasing the optimization budget (e.g., 1000 trials per feature), and
best configuration settings (e.g.\ transform types, ML model) can vary across
datasets. The greedy iterative search may converge to local optima and miss globally
optimal feature combinations that would require joint optimization. Searching over
segment boundaries expands the search space, which is both a source of
flexibility and of computational and overfitting risk. The formulation assumes
fixed-length, complete time series; streaming or incomplete monitoring (e.g., observation window not yet finished) is not addressed. Our
incremental feature analysis (Section~\ref{sec:feature_analysis}) reveals that the
feature search can overfit to the optimization split, particularly for cross-subject
tasks like PAMAP2 where the internal validation set shares subjects with training but
the test set contains an entirely unseen subject (validation--test gap
${\sim}26\%$). The choice of baselines and datasets carries a cherry-picking risk that
may limit the transferability of our conclusions to other task and method selections.

Our baseline comparison does not include several state-of-the-art methods such as HIVE-COTEv2
\cite{middlehurst2021hivecote2}, InceptionTime \cite{fawaz2020inceptiontime}, or LSTM-based architectures \cite{ordonez2016deep_activity}, with strong performance on
standard benchmarks \cite{middlehurst2024bakeoff}. We do not
evaluate on the UCR archive \cite{dau2019ucr}. These choices are deliberate: 
we focus on a flexible framework for diverse real-world biomedical datasets with clinically meaningful endpoints and large sample sizes, 
where interpretability and domain alignment matter as much as accuracy.
Many UCR datasets are small and lack the clinical context to assess whether discovered features are meaningful.

\subsection{Future Directions}

Future work includes learned or domain-specific signal transformations (e.g.,
wavelets or task-specific filters) that could capture patterns beyond the default
library, more efficient search strategies (e.g., early stopping, warm-starting from
related tasks) to reduce computational cost, and additional search for signal
transformation functions or segment summary statistics for better discriminative
power.

The bottom-up philosophy underlying SublimeX---starting minimal and iteratively adding complexity guided by task
performance---could apply to other model families: e.g., iteratively adding layers or components to a neural network until no further improvement, yielding
compact task-adapted architectures. Such extensions would generalize the paradigm beyond feature engineering to model construction.
The framework is actively maintained; we welcome contributions to extend it to new domains and modalities. 

\section{Conclusion}
We introduced SublimeX, an interpretable feature extraction framework built around three core principles: minimal feature sets discovered
through iterative optimization, interpretable features with explicit descriptions, and a flexible framework enabling domain-specific customization.

Across seven biomedical datasets, we demonstrated that SublimeX achieves performance comparable to or exceeding
state-of-the-art methods while providing transparent, human-readable
features with reduced dimensionality. SublimeX achieves the best
performance on AZT1D glucose forecasting (31.11 RMSE) and ties with MiniRocket
for the best AUC on REMC gene expression prediction ($0.931 \pm 0.023$ across
56 cell lines), with competitive performance on MIMIC-IV ($0.831 \pm 0.020$; tsfresh leads at $0.839$), MITBIH and Emotions. 
For highly specific tasks like PAMAP2 activity recognition or SVD voice pathology detection, SublimeX still struggles. 
Domain-specific SHAP analysis on two contrasting tasks---gene
expression prediction from histone modifications and glucose
forecasting from continuous monitoring data---reveals that
discovered features align with established clinical and
biological knowledge in both domains, demonstrating autonomous
discovery of physiologically meaningful characteristics.

SublimeX addresses a critical need in biomedical machine learning: methods that
are not only accurate but also transparent enough for clinical validation and scientific insight. 

\section*{Declaration of Competing Interest}

The authors declare that they have no known competing financial interests or personal relationships that could have appeared to influence the work reported in this paper.

\section*{Funding}
This work was supported by the Interdisciplinary Center for Clinical Research (IZKF) of the Medical Faculty of RWTH Aachen University.

\section*{CRediT Authorship Contribution Statement}

\textbf{J.C. Wolber}: Conceptualization, Methodology, Software, Validation,
Formal analysis, Investigation, Data curation, Writing -- original draft,
Writing -- review \& editing, Visualization.
\textbf{N.B. Paul}: Methodology, Software, Validation, Formal analysis,
Writing -- review \& editing.
\textbf{M. E. Samadi}: Methodology, Software, Validation, Formal analysis,
Writing -- review \& editing.
\textbf{J. Sellin}: Writing -- review \& editing, Supervision.
\textbf{M. Muecke}: Writing -- review \& editing, Supervision, Project
administration.
\textbf{A. Schuppert}: Writing -- review \& editing, Project administration, Supervision.
\section*{Data Availability}

The SublimeX Python package is available via PyPI.
Source code and documentation are available at \url{https://github.com/Prgrmmrjns/sublimex} and \url{https://pypi.org/project/sublimex/}.
The manuscript repository containing evaluation scripts is available at \url{https://github.com/Prgrmmrjns/patternExtraction-manuscript}.

Dataset descriptions, sources, and access are in Section~\ref{sec:dataset_processing} of the Supplementary Material.

\bibliographystyle{unsrt}
\bibliography{bibliography}



\section{Supplementary Material}

\subsection{Baseline Implementations}

We compare SublimeX against five time series methods, ranging from feature engineering to deep learning methods.
For feature-based methods (SublimeX, tsfresh, catch22, MiniRocket, RDST), we train a LightGBM classifier on the extracted features. CNN is trained end-to-end.

\subsubsection{tsfresh}
\textbf{tsfresh} \cite{christ2018tsfresh} extracts a comprehensive library of statistical and spectral characteristics from time series data.
We use \texttt{MinimalFCParameters} to balance feature coverage and computational cost, yielding 8 features per univariate series.
For multivariate time series, we extract features independently for each channel and concatenate them into a single feature vector.

\subsubsection{catch22}
\textbf{catch22} \cite{lubba2019catch22} provides a canonical set of 22 computationally efficient time series features, selected from the larger hctsa
library for their diversity and discriminative power.
For multivariate data, we compute all 22 features per channel and concatenate them, resulting in $22 \times n\_channels$ features.

\subsubsection{MiniRocket}
\textbf{MiniRocket} \cite{dempster2021minirocket} transforms time series via approximately 10,000 random convolutional kernels with fixed (non-learned) weights.
For each kernel, it extracts the proportion of positive values (PPV) after convolution, yielding a high-dimensional feature representation.
We use \texttt{MiniRocketMultivariate} from the sktime library, which natively handles multivariate inputs by applying kernels across all channels.
This method represents the ``accuracy-at-all-costs'' end of the spectrum, offering high performance but limited interpretability.

\subsubsection{RDST}
\textbf{RDST} \cite{guillaume2022random} (Random Dilated Shapelet Transform) is a fast and accurate shapelet-based approach that introduces dilation to capture 
shapelets at multiple scales with improved computational efficiency.
We configure RDST with \texttt{max\_shapelets=1000}. The transform generates distance-based features measuring how similar each time series is to the discovered shapelets.

\subsubsection{CNN}
\textbf{CNN}: A 1D Convolutional Neural Network trained end-to-end, representing the deep learning baseline.
We implement a 1D CNN using PyTorch with three convolutional layers (filters: 64, 128, 256) each containing convolution, batch normalization, ReLU
activation, max pooling, and dropout (0.2, 0.2, 0.3). The architecture uses kernel sizes of 7, 5, and 3 for the three layers
respectively, with appropriate padding to maintain temporal resolution. Adaptive average pooling aggregates features across the temporal dimension,
followed by fully connected layers (256$\rightarrow$128$\rightarrow$64$\rightarrow$output) with ReLU activations, batch normalization, and dropout (0.4, 0.3) for regularization.

We train with Adam optimizer (lr=$10^{-3}$, weight decay=$10^{-4}$), early stopping (patience=10), and learning rate reduction on plateau (factor=0.5, patience=3).
Input data is z-score normalized using training set statistics. The model automatically handles multivariate inputs by adjusting the number of input channels 
based on the dataset structure. Training uses 300 epochs with batch size 128 and 80/20 train/validation split for early stopping.

\subsection{Dataset Processing Details}
\label{sec:dataset_processing}

\subsubsection{AZT1D Preprocessing}

The AZT1D dataset is available via Zenodo (\url{https://zenodo.org/records/15094234}).
For this dataset, we implemented knowledge-driven response curve modeling based on established pharmacokinetic and glycemic response literature.
For insulin, we use a gamma-like function with 15-minute onset, 60-minute peak,
and 4-hour duration; for carbohydrates, we use 10-minute onset, 45-minute peak, and 3-hour duration.
These response curves model the temporal dynamics of physiological effects,
with insulin effects peaking around 60 minutes and persisting for up to 4
hours, while carbohydrate effects peak around 45 minutes and persist for up to
3 hours.

At each timepoint, we compute cumulative active effects by convolving all historical insulin doses and carbohydrate intakes with their respective
response curves. This creates interpretable time series that capture overlapping physiological
effects, enabling the model to understand how multiple insulin doses and meal intakes interact over time.

The preprocessing pipeline extracts 24 timepoints (2 hours of history at
5-minute intervals) for three physiological signals: CGM values, active insulin effects, and active carbohydrate effects.
The regression target is future glucose change $g_{t+h} - g_t$ where $h = 60$ minutes (12 timepoints ahead).
We also include the initial CGM value (CGM\_0) as a baseline feature to provide the current glucose context.

\begin{figure}[H]
  \centering
  \includegraphics[width=\textwidth,keepaspectratio]
  {images/supplementary_data/azt1d_insulin_carb_visualization.png}
  \caption{AZT1D preprocessing visualization showing insulin and carbohydrate response modeling.
  Top panel: Real-time CGM with overlapping insulin and carbohydrate effects over time for Subject 1.
  Blue line shows CGM values with target range (70-180 mg/dL) highlighted in green.
  The dark green shaded area represents active insulin effect from overlapping doses (circles with crosses).
  The orange shaded area represents active carbohydrate effect from overlapping intakes (circles with crosses).
  Bottom panel: Physiological response curves showing normalized insulin (15-min onset, 60-min peak, 240-min duration) and carbohydrate (10-min onset, 45-min peak, 180-min duration) effects over time.}
\end{figure}

\subsubsection{MITBIH Preprocessing}

The MIT-BIH Arrhythmia Database is available via PhysioNet (\url{https://physionet.org/content/mitdb/1.0.0/}).
The database contains annotated ECG recordings from 48 patients.
We extract individual heartbeats using a fixed window size around annotated R-peak positions (100 samples before and 100 samples after, totaling 200 samples per beat).
When the window extends beyond signal boundaries, we pad with zeros to maintain consistent segment length.
Each beat is classified into one of five arrhythmia types: Normal (N), Left bundle branch block (L), Right bundle branch block (R), Atrial premature (A), and Premature ventricular contraction (V).

\subsubsection{Emotions Preprocessing}

The Emotions EEG dataset is available via Kaggle (\url{https://www.kaggle.com/datasets/birdy654/eeg-brainwave-dataset-feeling-emotions}).
The dataset contains brainwave recordings collected while participants experienced different emotional states (positive, neutral, negative).
We use raw time-domain EEG signals from 5 channels (AF3, F7, F3, FC5, T7) with 254 time points per channel, without any frequency-domain transformations or filtering.

\subsubsection{REMC Preprocessing}

The Roadmap Epigenomics Consortium data is available at \url{http://www.roadmapepigenomics.org/}.
Following our prior work on PatternChrome \cite{paul2025prediction}, we predict binary gene
expression (high/low) from histone modification ChIP-seq profiles from the Roadmap Epigenomics Consortium (all 56 cell lines, 19,795 genes).
We convert RefSeq to GENCODE identifiers, retaining 18,421 genes.
For each gene, we extract five histone modification signals (H3K4me1, H3K4me3, H3K36me3, H3K9me3, H3K27me3) from 5000 bp upstream and downstream of the TSS,
divided into 100 bp bins (100 bins per mark). Gene expression is binarized at the median level.

\subsubsection{MIMIC-IV Preprocessing}

The MIMIC-IV database is available via PhysioNet
(\url{https://physionet.org/content/mimiciv/}; requires credentialing)
\cite{Johnson2023MIMICIV}.
We predict mortality from ICU patients on mechanical ventilation using this database.
The observation window comprises the first 24 hours of mechanical ventilation.
Mortality is defined as death within 6 days after the observation window.

The dataset includes static features, aggregated clinical measurements, and multivariate time
series. Static features (21 total) include demographics (age, gender), comorbidities (diabetes,
hypertension, immunosuppression, renal failure, obesity, heart failure, liver disease, cancer),
BMI, and admission service categories. We exclude diagnostic flags (ARDS, Sepsis-3, SIRS,
anemia) from initial features to prevent data leakage, as these flags are derived from clinical
measurements already captured by the time series and aggregated features.
Aggregated features (13 total) are clinical measurements with low temporal
variation, averaged over the 24-hour window: laboratory values (creatinine,
platelets, WBC, INR, lactate, bilirubin), blood gases (PaCO2, pH), GCS verbal
and motor components, and vasopressor doses (norepinephrine, epinephrine, dobutamine).
Time series features (12 channels $\times$ 24 hours) capture measurements with
high temporal variation: vital signs (heart rate, respiratory rate, SpO2,
blood pressure, temperature), ventilator settings (FiO2, PEEP), PaO2, GCS eye opening, and urine output.
Subjects with fewer than 24 hours of data are excluded. Missing values in time series are imputed using forward-fill followed by
backward-fill along the temporal axis, with remaining gaps filled with zero.


\subsubsection{PAMAP2 Preprocessing}

The PAMAP2 dataset is available via the UCI ML Repository
(\url{https://archive.ics.uci.edu/dataset/231/pamap2+physical+activity+monitoring}).
The dataset contains recordings from 9 subjects performing 18 physical
activities while wearing 3 IMUs (wrist, chest, ankle) and a heart rate monitor.
Subject 109 is excluded from training and evaluation as they only performed one activity
(rope jumping), making LOSO evaluation invalid for that fold (8 subjects used).
Following the preprocessing protocol from Tello et al. \cite{tello2025contrasting}, we
use 27 sensor channels: 3D acceleration ($\pm$16g scale), 3D gyroscope, and
3D magnetometer from each of the three IMUs (9 channels $\times$ 3 IMUs).
We exclude temperature, heart rate, 3D acceleration ($\pm$6g scale), and 4D orientation channels.
We create 512-sample windows (5.12 seconds at 100Hz) with a 100-sample shift
(1 second), filtering transient activities (ID 0) to obtain 12 activity
classes: lying, sitting, standing, walking, running, cycling, Nordic walking,
watching TV, computer work, ascending/descending stairs, and vacuum cleaning.
Missing values are interpolated linearly and remaining NaNs filled with column means.

\subsubsection{SVD Preprocessing}

The Saarbr\"ucken Voice Database (SVD) \cite{svd_database} is available at \url{https://stimmdb.coli.uni-saarland.de}.
It is a publicly available collection of voice recordings from over 2{,}000
subjects, including healthy controls and patients with 71 different voice pathologies \cite{martinez2012voice}.
We use recordings of the phrase ``Guten Morgen, wie geht es Ihnen?'' (Good morning, how are you?).
All pathology types are included without filtering, representing a realistic hospital triage scenario.
Audio is truncated to 3 seconds maximum and resampled to 1{,}000 samples.
The final dataset contains 1{,}988 recordings: 634 healthy controls (31.9\%) and 1{,}354 pathological cases (68.1\%).

\end{document}
